%% kapa: exam=1 round=33 year=2022 period=2 form=A subject=accounting subject_ko=회계학 seq=4 number=44
%% source: Q-Net 감정평가사 기출문제 (cst00309) artlSeq=5209997 file=raw/1008/5209997/2022년도 제33회 감정평가사 1차시험 2교시(감정평가 관계법규, 회계학).pdf
%% conversion: kordoc 4.24.51 -> tools/gen1.py (auto); review=proofread-sonnet; flags=none
\begin{question}{44}
(주)감평은 20×1년 초 임대목적으로 건물(취득원가 ￦1,000, 내용연수 10년, 잔존가치 ￦0, 정액법 감가상각)을 취득하여 이를 투자부동산으로 분류하였다. 20×1년 말 건물의 공정가치가 ￦930일 때 (A)공정가치모형과 (B)원가모형을 각각 적용할 경우 (주)감평의 20×1년도 당기순이익에 미치는 영향은? (단, 해당 건물은 매각예정으로 분류되어 있지 않다.)
\begin{choices}
  \item (A) ￦70 감소 \quad (B) ￦100 감소
  \item (A) ￦70 감소 \quad (B) ￦70 감소
  \item (A) ￦30 감소 \quad (B) ￦100 감소
  \item (A) ￦30 증가 \quad (B) ￦70 감소
  \item (A) ￦30 증가 \quad (B) ￦30 증가
\end{choices}
\end{question}
\answer{1}
